\documentclass[../../main.tex]{subfiles}
\begin{document}

{\bf [解]}

漸化式
\begin{align}
  \begin{cases}
    u_1 = 2 \\
    u_2 = a^2+2 \\
    u_{n+2} = au_{n}-u_{n+1}
  \end{cases} \label{eq:1}
\end{align}
について考える．
初めの4項は
\begin{align}
  u_1=2，\quad u_2=a^2+2，\quad u_3=2a-a^2-2，\quad
  u_4=a^3+a^2+2 \label{eq:2}
\end{align}
となる．以下，合同式はすべて法4で考える
\cref{eq:2}において$a\equiv1$のとき$u_4\equiv0$となるので不適である．
従って$a\equiv0,2,-1$が必要である．以下各ケースについて個別に考える．

\subparagraph{$a\equiv0$のとき．}

すべての$n$について$u_n\not\equiv0$となることを数学的帰納法で示す．
$n=1,2$では$u_1\equiv u_2\equiv2$である．
そこで$k\in\mathbb{N}$に対して$n=k,k+1$での成立を仮定すれば，漸化式\cref{eq:1}から
\begin{align}
  u_{k+2}\equiv au_k - u_{k+1} \equiv -u_{k+1}\not\equiv0
\end{align}
となり，$n=k+2$でも成立する．従って数学的帰納法によりすべての$n$で$u_n\not\equiv0$であり，$4$の倍数でない．

\subparagraph{$a\equiv2$のとき．}
すべての$n$について$u_n\equiv2$となることを数学的帰納法で示す．
$n=1,2$では\cref{eq:2}より成立する．
そこで$k\in\mathbb{N}$に対して$n=k,k+1$での成立を仮定すれば，漸化式\cref{eq:1}から
\begin{align}
  u_{k+2} \equiv au_k - u_{k+1} \equiv 2u_k-u_{k+1}\equiv 4-2 \equiv 2
\end{align}
となり，$n=k+2$でも成立する．従って数学的帰納法によりすべての$n$で$u_n\equiv2$である．よって全ての$u_n$は$4$の倍数でない．

\subparagraph{$a\equiv-1$のとき．}
$u_n$は法4の剰余類で$(2,-1,-1)$の周期数列となることを数学的帰納法で示す．
$n=1,2,3$では実際に$(u_1,u_2,u_3)\equiv(2,-1,-1)$である．
ある$k\geqq1$について，$u_{3k-2},u_{3k-1},u_{3k}$での成立を仮定すると，漸化式から
\begin{align}
  u_{3k+1}&\equiv au_{3k-1} - u_{3k}   \equiv-(-1)-(-1)\equiv2，\\
  u_{3k+2}&\equiv au_{3k}   - u_{3k+1} \equiv -2-(-1)\equiv-1，\\
  u_{3k+3}&\equiv au_{3k+1} - u_{3k+2} \equiv -(-1)-2\equiv-1
\end{align}
となる．従って数学的帰納法により$(2,-1,-1)$の周期数列となり，$4$の倍数となる項は存在しない．
\casesend

以上より，必要十分条件は
\begin{align}
  a\equiv-1,0,2\pmod{4}
\end{align}
である．

\newpage
\end{document}
